\BeleskeID{09_3-s063}
% element: 09_3-s063:e04
% slide-id: 09_3-s063
Za N tranzistor razmotriti oba moguća izbora sorsa. Ako bi sors bio na strani 1 uz struju iz mase kroz $R_L$, drejn bi bio negativan u odnosu na sors, što nije saglasno tom izboru. Ako bi sors bio na strani 2 uz struju ka masi, njegov potencijal bi bio veći od potencijala drejna, što ponovo nije saglasno izboru nižeg potencijala kao sorsa.

U ravnoteži je zato struja nula i oba kraja su na istom potencijalu. N tranzistor ima uslove za provođenje i radi u omskoj oblasti sa malom dinamičkom otpornošću iako je jednosmerna struja nula. Analogno razmišljanje važi i za P tranzistor. Izlaz ostaje $0\,\mathrm V$.
